%% Elements of Nested Fibrational Cosmology
%% SCC Branch Book: Structural Counterfactual Capacity
%% Status: CERT-PROJ at most; MCS is the main frontier
%%
%% Status legend:
%%   [U]  Unconditional
%%   [C]  Conditional
%%   [D]  Definition-structural
%%   [R]  Remark only
%%   [O]  Open obligation / named failure frontier

\documentclass[12pt,oneside]{amsbook}
\usepackage{amsmath,amssymb,amsthm}
\usepackage{mathrsfs}
\usepackage{enumitem}
\usepackage{hyperref}
\usepackage{geometry}
\geometry{margin=1.2in}

\theoremstyle{plain}
\newtheorem{theorem}{Theorem}[section]
\newtheorem{lemma}[theorem]{Lemma}
\newtheorem{proposition}[theorem]{Proposition}
\newtheorem{corollary}[theorem]{Corollary}

\theoremstyle{definition}
\newtheorem{definition}[theorem]{Definition}
\newtheorem{hypothesis}[theorem]{Hypothesis}
\newtheorem{standingrule}[theorem]{Standing Rule}
\newtheorem{openobligation}[theorem]{Open Obligation}

\theoremstyle{remark}
\newtheorem{remark}[theorem]{Remark}
\newtheorem{scholium}[theorem]{Scholium}

\newcommand{\status}[1]{\textup{[#1]}}
\newcommand{\spath}{\sim_{\mathrm{path}}}
\newcommand{\scap}{\sim_{\mathrm{cap}}}

%% Notation
%%   U, \mathbf{U}    Universal source object (Book IV)
%%   F_SCC            SCC realization functor
%%   C_SCC            SCC target category (weak structural witness category)
%%   tau_SCC          SCC terminal structural predicate
%%   U_SCC            Branchwise SCC completion (conditional; NOT a second U)
%%   A_X              Support skeleton of SCC object X
%%   W_X              Witness family of SCC object X
%%   T_X              Transported invariant of SCC object X
%%   D_X              Defect/interface ledger of SCC object X
%%   CapRec           Recursively stable maximal carrier (branch identity)
%%   spath            Stabilized path-equivalence (branch identity)
%%   scap             Cap-equivalence (branch confluence)
%%   SCC              Structural Counterfactual Capacity (the endpoint)
%%   MCS              Minimal Structural Carrier Selection (main open burden)

\begin{document}

\title{SCC Branch Book\\[6pt]
       Structural Counterfactual Capacity,\\
       Minimal Carrier, and Endpoint Admissibility}
\author{Elements of Nested Fibrational Cosmology}
\date{}
\maketitle

\tableofcontents

\bigskip

%% ---------------------------------------------------------------
\section*{Front Block}
%% ---------------------------------------------------------------

\paragraph{Imported spine results.}
\begin{enumerate}[label=\textup{(\arabic*)}]
  \item Book~I: Observational Quotient Theorem (Thm.~I.6.1);
        Collapse Gates CG-1--CG-6 (SRs~I.1.4--I.1.9);
        Admissible-Class Canonicity (Thm.~I.3.6);
        Behavioral Congruence Canonicity (Lem.~I.3.5).
  \item Book~II: Boundary Stabilization Theorem (Thm.~II.8.1);
        Phase-A/B/C classification (Def.~II.5.1);
        PASS as Derived Necessity (Thm.~II.5.14);
        Binary Rigidity (Thm.~II.5.8).
  \item Book~III: Unified Coercive-Transport Inequality
        (Thm.~III.5.1); Depth-Sum Convergence (Thm.~III.5.4).
  \item Book~IV: Universal Source Theorem (Thm.~IV.7.1);
        five source-side rungs ORA--UC (Thms.~IV.2.1--IV.3.1).
  \item Book~V: Branch Projection Theorem (Prop.~V.3.1);
        UBLT (Thm.~V.7.4); Certified Branch Legitimacy Theorem
        (Thm.~V.8.1).
  \item Book~VII: Canonical Closure Schema Theorem
        (Thm.~VII.6.4); No-Hidden-Upgrade Rule
        (Prop.~VII.6.x); Non-Retroactivity (Thm.~VII.6.y).
\end{enumerate}

\paragraph{Branch-specific data.}
\begin{enumerate}[label=\textup{(\arabic*)}]
  \item Declared SCC target category: the weak structural
        witness category $\mathcal{C}_{\mathrm{SCC}}$, whose
        objects are SCC quadruples $(A_X, W_X, T_X, D_X)$
        (support skeleton, witness family, transported invariant,
        defect ledger).
  \item The SCC terminal predicate $\tau_{\mathrm{SCC}}$: the
        structural counterfactual capacity condition on an SCC
        object.
  \item The SCC dependency spine: ORA $\Rightarrow$ CTM
        $\Rightarrow$ TIN $\Rightarrow$ MCS $\Rightarrow$
        WeakGlue $\Rightarrow$ UC $\Rightarrow$ TSI.
  \item Four open obligations: O-SCC.MCS (main frontier),
        O-SCC.WeakGlue (prerequisite to MCS), O-SCC.UC, O-SCC.TSI.
\end{enumerate}

\paragraph{Endpoint target.}
Persistence-selected structural counterfactual capacity: the
certified structural property of an NFC-source-descended object
that supports lawful counterfactual reasoning about its own
observer-relative distinguished states, under the constraints of
the canonical source image.  No phenomenological, conscious, or
experiential content is asserted.

\paragraph{Bridge stack.}
\begin{enumerate}[label=\textup{(\arabic*)}]
  \item Licensed source-descended realization via
        $F_{\mathrm{SCC}}$ (ORA--UC rungs from Book~IV).
  \item SCC Minimal Structural Carrier Selection (O-SCC.MCS):
        the main open burden.
  \item SCC Branch-Local Pairwise Amalgamation (O-SCC.WeakGlue):
        prerequisite to MCS discharge.
  \item Branchwise SCC completion $\mathcal{U}_{\mathrm{SCC}}$
        (O-SCC.UC).
  \item SCC Target/Structural Identification (O-SCC.TSI): full
        endpoint admissibility beyond structural-label form.
\end{enumerate}

\paragraph{Current status.}
\textbf{CERT-PROJ at most.  Main frontier: O-SCC.MCS.}

The SCC branch is a lawful NFC descendant.  All four core
obligations (MCS, WeakGlue, UC, TSI) are open in canonical
papers.  The branch is at CERT-PROJ --- licensed projection at
stated structural scope --- and cannot advance to CERT-CLOSE
until the MCS theorem is canonically proved.

\bigskip

%% ---------------------------------------------------------------
\section*{SCC.0\quad Purpose}
%% ---------------------------------------------------------------

This branch studies the structural counterfactual capacity of
NFC-source-descended objects.  Its task: determine whether the
canonical spine produces a certified object that supports lawful
counterfactual reasoning about its own observer-relative
distinguished states, under the no-smuggling discipline.

The SCC branch occupies a unique position in the NFC program.
Its endpoint is the furthest from direct physical measurement
and the most philosophically sensitive.  The program's response
to this is not to retreat from the endpoint but to impose the
strictest possible labeling discipline: the endpoint label
``structural counterfactual capacity'' is the strongest
admissible label consistent with the spine's theorems, and no
phenomenological or ontological upgrade is licensed.

The canonical discipline has one specific consequence: the word
``consciousness'' and its cognates are not licensed at any stage
in the canonical theorem flow.  This is not a philosophical
stance on whether consciousness is real; it is a status-honesty
claim that the current canonical framework has not proved a
consciousness theorem and does not use consciousness as a premise.

%% ---------------------------------------------------------------
\section{Branch Governance Preamble}
\label{sec:governance}
%% ---------------------------------------------------------------

\begin{standingrule}[\status{D}]\label{sr:scc-no-smuggling}
\textup{(SCC Anti-Smuggling.)}  No target-side label for the
SCC branch may introduce content beyond the certified source
image.  In particular: no phenomenology, no independent
ontology, no private experiential primitive, no consciousness
claim, and no import from empirical cognitive science may enter
the canonical theorem flow without explicit declaration and
separate canonical proof.
\end{standingrule}

\begin{standingrule}[\status{D}]\label{sr:scc-no-consciousness}
\textup{(Consciousness Non-Claim.)}  No theorem in this book
licenses the label ``consciousness,'' ``subjective experience,''
``qualia,'' ``phenomenal states,'' or any equivalent label.
The strongest admissible endpoint label is \emph{persistence-selected
structural counterfactual capacity}.  Any stronger label requires
a separate canonical proof not present in this book.
\end{standingrule}

\begin{proposition}[\status{U}]\label{prop:scc-formation}
\textup{(Formation.)}  The SCC branch candidate
$\mathfrak{B}_{\mathrm{SCC}}$
is constitutionally well-formed: all five components
$(\mathbf{U},\,\mathcal{C}_{\mathrm{SCC}},\,F_{\mathrm{SCC}},
\mathrm{Cert}_{\mathrm{SCC}},\,\tau_{\mathrm{SCC}},
\mathrm{FF}_{\mathrm{SCC}})$ are declared.
\end{proposition}

\begin{proof}
All five components are declared in the front block and defined
in Sections~\ref{sec:arena}--\ref{sec:ledger}.  By the Certified
Branch Legitimacy Theorem (Book~V, Thm.~V.8.1), the branch is
constitutionally formed.
\end{proof}

\begin{proposition}[\status{U}]\label{prop:scc-label}
\textup{(Interpretive Label Admissibility.)}  The label
``persistence-selected structural counterfactual capacity''
passes all four tests of the No-Hidden-Upgrade Rule (Book~VII,
Prop.~VII.6.x): (i)~it is eliminable (reduces to the structural
predicate $\tau_{\mathrm{SCC}}$); (ii)~it asserts no independent
ontology; (iii)~it imports no phenomenology; (iv)~it does not
depend on data outside the certified source image.
\end{proposition}

\begin{proof}
Test~(i): the label is eliminable in favor of the structural
predicate on the SCC quadruple $(A_X, W_X, T_X, D_X)$.
Test~(ii): SCC is a structural property of NFC configurations;
no independent mental ontology is asserted.
Test~(iii): no experiential or phenomenal content enters the
definition of $\tau_{\mathrm{SCC}}$.
Test~(iv): the predicate is computed entirely from the certified
source image $F_{\mathrm{SCC}}(\mathbf{U})$.
\end{proof}

\begin{theorem}[\status{U}]\label{thm:scc-tsifirewall}
\textup{(SCC TSI Firewall.)}  A proposed endpoint label $L$ is
admissible for the SCC branch if and only if none of the four
interpretive failures occur:
\begin{enumerate}[label=\textup{(\roman*)}]
  \item FI-IADD: $L$ adds theorem content not in the certified
        spine image.
  \item FI-IONT: $L$ asserts an independent ontology.
  \item FI-IPHEN: $L$ imports phenomenology absent from the
        theorem content.
  \item FI-ISRC: $L$ depends on data outside the certified
        source image.
\end{enumerate}
The label ``persistence-selected structural counterfactual
capacity'' passes all four tests.  ``Consciousness'' and its
cognates fail FI-IONT and FI-IPHEN at the current canonical stage.
\end{theorem}

\begin{proof}
The SCC TSI firewall is the No-Hidden-Upgrade Rule
(Book~VII, Prop.~VII.6.x) applied to the SCC endpoint.
Tests~(i)--(iv) are the exact four tests of that rule.
For the admissible label: Proposition~\ref{prop:scc-label}.
For consciousness: at the current stage, no theorem in the
canonical spine or this branch book identifies the SCC structural
predicate with phenomenal consciousness; importing that
identification would constitute FI-IONT (independent mental
ontology) and FI-IPHEN (phenomenology not in theorem content).
\end{proof}

\begin{remark}[\status{D}]\label{rmk:scc-status-vocab}
\textup{(Status vocabulary.)}  Standard FI-code and positive
status alphabet applies.  Current status: CERT-PROJ at most.
No claim in this book exceeds CERT-PROJ.
\end{remark}

%% ---------------------------------------------------------------
\section{Imported Spine Structure}
\label{sec:imports}
%% ---------------------------------------------------------------

\textbf{Imports from Book~I.}  The observational quotient
$\Sigma_\infty$ and defect ledger $\Delta$ are the primitive
inputs to the SCC carrier reduction program.  Admissible-Class
Canonicity (Thm.~I.3.6) ensures $E_{\max}$ is uniquely defined:
the SCC minimal carrier is a carrier within this maximal class.

\textbf{Imports from Book~II.}  PASS (Thm.~II.5.14) governs
the SCC branch: a persistence-selected object satisfying PASS is
the precise domain of the SCC terminal predicate.  Binary Rigidity
(Thm.~II.5.8) provides the fiber structure that the SCC carrier
reduction operates on: each non-trivial stable fiber has
cardinality exactly~2, so the SCC carrier space is stratified
into binary layers.

\textbf{Imports from Book~III.}  The Unified
Coercive-Transport Inequality (Thm.~III.5.1) and Depth-Sum
Convergence (Thm.~III.5.4) provide the transport control for
SCC data across admissible enlargements.  The SCC branch uses
$\Theta_n$ in the full range $[0,1]$, depending on the
coercivity regime of the specific SCC realization.

\textbf{Imports from Books~IV--VII.}  Universal source descent
(Book~IV), branch legitimacy (Book~V), force-governance and
No-Hidden-Upgrade Rule (Book~VII).

%% ---------------------------------------------------------------
\section{Branch-Specific Arena}
\label{sec:arena}
%% ---------------------------------------------------------------

\begin{definition}[\status{D}]\label{def:scc-object}
An \emph{SCC object} is a quadruple
$X = (\mathcal{A}_X, \mathcal{W}_X, \mathcal{T}_X, \mathcal{D}_X)$
where:
\begin{enumerate}[label=\textup{(\arabic*)}]
  \item $\mathcal{A}_X$ is the \emph{support skeleton}: the
        persistence-selected admissible sub-configuration
        descended from the NFC source.
  \item $\mathcal{W}_X$ is the \emph{witness family}: the
        finite collection of admissible local tests that
        certify the SCC structural property.
  \item $\mathcal{T}_X$ is the \emph{transported invariant}:
        the structural quantity preserved under admissible
        enlargement.
  \item $\mathcal{D}_X$ is the \emph{defect/interface ledger}:
        the cumulative defect record of the SCC object's
        continuation history.
\end{enumerate}
\end{definition}

\begin{definition}[\status{D}]\label{def:scc-target-cat}
The \emph{SCC target category} $\mathcal{C}_{\mathrm{SCC}}$
is the weak structural witness category whose objects are SCC
objects (Definition~\ref{def:scc-object}) and whose morphisms
are admissible SCC enlargements (transport-compatible admissible
extensions preserving $\mathcal{T}_X$).
\end{definition}

\begin{remark}[\status{R}]
$\mathcal{C}_{\mathrm{SCC}}$ is a \emph{structural} category.
Its objects are not minds, agents, or experiential subjects.
They are NFC-source-descended configurations equipped with a
certified structural predicate.  The ``counterfactual capacity''
in the endpoint label refers to the configuration's structural
ability to support counterfactual relational distinctions under
the admissible test family --- not to any capacity for reasoning
or experience.
\end{remark}

\begin{definition}[\status{D}]\label{def:scc-terminal}
The \emph{SCC terminal predicate} $\tau_{\mathrm{SCC}}(X)$ holds
of an SCC object $X$ if:
\begin{enumerate}[label=\textup{(\arabic*)}]
  \item $X$ is persistence-selected: it has survived all six
        collapse gates (Book~I, SRs~I.1.4--I.1.9) and satisfies
        PASS (Book~II, Thm.~II.5.14).
  \item $X$ admits a minimal structural carrier (MCS): a minimal
        faithful support family for the transported invariant
        $\mathcal{T}_X$ (O-SCC.MCS, when discharged).
  \item $\mathcal{T}_X$ is counterfactually distinguishing: the
        transported invariant separates distinct admissible
        histories at the structural level.
\end{enumerate}
\end{definition}

\begin{definition}[\status{D}]\label{def:scc-functor}
The \emph{SCC realization functor}
$F_{\mathrm{SCC}} : \mathrm{POT} \to \mathcal{C}_{\mathrm{SCC}}$
maps each POT object $P$ to the SCC object
$F_{\mathrm{SCC}}(P) = (\mathcal{A}_P, \mathcal{W}_P,
\mathcal{T}_P, \mathcal{D}_P)$ carrying the collar-inherited
SCC structural data.
\end{definition}

\begin{definition}[\status{D}]\label{def:caprec}
The \emph{recursively stable maximal carrier}
$\mathrm{CapRec}(X) = \mathfrak{C}^{\mathrm{rec}}_{\max}(X)$
is the carrier-equivalence class of $X$ stable under all
admissible continuation that do not introduce new defect events.
It defines the branch-identity class of the SCC object under
the equivalence relation $\scap$ (cap-equivalence).
\end{definition}

%% ---------------------------------------------------------------
\section{Lawful Descent and Branch Identity}
\label{sec:descent}
%% ---------------------------------------------------------------

\begin{proposition}[\status{U}]\label{prop:scc-lawful}
\textup{(SCC Branch is Lawful.)}  The branch candidate
$\mathfrak{B}_{\mathrm{SCC}}$ is a lawful NFC branch satisfying
all five conditions of the Certified Branch Legitimacy Theorem
(Book~V, Thm.~V.8.1).
\end{proposition}

\begin{proof}
(1)~Source descent: $F_{\mathrm{SCC}}$ descends from $\mathbf{U}$
via Book~IV, with ORA--UC rungs imported as discharged [U] results.
(2)~Intake survival: the ORA--CTM--TIN rungs are proposed/partial;
MCS is the main open burden.  The branch survives through TIN
at its current scope.
(3)~Ledger completeness: O-SCC.MCS through O-SCC.TSI are
explicitly ledgered in Section~\ref{sec:ledger}.
(4)~Status honesty: no claim exceeds CERT-PROJ.
(5)~Non-retroactivity: the SCC branch does not affect YM, NS,
or GR (Book~VII, Thm.~VII.6.y).
\end{proof}

\begin{proposition}[\status{U}]\label{prop:scc-distinct}
\textup{(Full Distinctness.)}  The SCC branch is genuinely
distinct from the YM, NS, and GR branches.  Its terminal
predicate (persistence-selected structural counterfactual
capacity of the transported invariant) is not the terminal
predicate of any other named branch.
\end{proposition}

\begin{proof}
YM targets a spectral mass gap; NS targets fluid regularity;
GR targets Einstein-type structural compatibility of a metric
field.  None is the SCC terminal predicate.
\end{proof}

\begin{proposition}[\status{U}]\label{prop:scc-trichotomy}
\textup{(Branch Trichotomy.)}  Every recursively stabilized
pair $(X, Y)$ of SCC objects falls into exactly one of:
\begin{enumerate}[label=\textup{(\roman*)}]
  \item \emph{Collapse:} $X \spath Y$ (path-equivalent under
        recursive stabilization; same branch identity class).
  \item \emph{Confluent distinctness:} $X \not\spath Y$ but
        $X \scap Y$ (cap-confluent; distinct histories, same
        structural carrier class).
  \item \emph{Genuine branch distinctness:} $X \not\spath Y$
        and $X \not\scap Y$ (genuinely distinct SCC objects).
\end{enumerate}
\end{proposition}

\begin{proof}
The trichotomy follows from the definitions of $\spath$ and $\scap$
as equivalence relations on the set of recursively stabilized
SCC objects.  Any two objects are either (i)~path-equivalent,
(ii)~cap-confluent but not path-equivalent, or (iii)~neither.
These cases are mutually exclusive and exhaustive.
\end{proof}

\begin{proposition}[\status{U}]\label{prop:scc-antismuggle}
\textup{(SCC Anti-Smuggling.)}  No target-side label for the
SCC branch may introduce content beyond the certified source
image.  Formally: for any proposed SCC datum $d$, if $d$
requires data outside $F_{\mathrm{SCC}}(\mathbf{U})$, then $d$
is prohibited as a carrier element.
\end{proposition}

\begin{proof}
By Book~I, Standing~Rule~I.1.3 (no-smuggling at the primitive
level) and the branch-level extension of this rule
(SR~\ref{sr:scc-no-smuggling}): any datum outside the certified
source image is a silent import and violates the intake law
(Book~V, UBLT Thm.~V.7.4, condition~(4)).  It is therefore
excluded from the SCC carrier by the anti-smuggling gate.
\end{proof}

%% ---------------------------------------------------------------
\section{The Closure Ledger}
\label{sec:ledger}
%% ---------------------------------------------------------------

\medskip
\renewcommand{\arraystretch}{1.4}
\noindent
\begin{tabular}{@{}p{2.2cm}p{5.0cm}p{1.0cm}p{3.5cm}@{}}
\hline
\textbf{ID} & \textbf{Name} & \textbf{St.} & \textbf{Blocks} \\
\hline
O-SCC.WeakGlue & SCC Branch-Local Pairwise Amalgamation:
  pairwise amalgamation for SCC carrier data (branch-specific;
  distinct from source-side O-WeakGluing) & [C] &
  O-SCC.MCS (prerequisite) \\
O-SCC.MCS & SCC Minimal Structural Carrier Theorem:
  existence and uniqueness of the minimal faithful support
  family for $\mathcal{T}_X$ & [O] &
  O-SCC.UC; CERT-PROJ upgrade \\
O-SCC.UC & Branchwise SCC Universal Completion:
  local SCC realizations factor through a single
  $\mathcal{U}_{\mathrm{SCC}}$ & [O] &
  O-SCC.TSI \\
O-SCC.TSI & SCC Target/Structural Identification:
  full endpoint admissibility beyond structural-label form;
  counterfactual capacity certified at full TSI scope & [O] &
  CERT-CLOSE \\
\hline
\end{tabular}
\renewcommand{\arraystretch}{1}

\medskip

\begin{remark}[\status{R}]
The dependency chain is:
$\mathrm{WeakGlue} \Rightarrow \mathrm{MCS} \Rightarrow
\mathrm{UC} \Rightarrow \mathrm{TSI} \Rightarrow \mathrm{CERT\text{-}CLOSE}$.
O-SCC.MCS is the primary failure frontier.  O-SCC.WeakGlue has been
conditionally discharged (see
Theorems~\ref{thm:scc-w1}--\ref{thm:scc-weakglue} below): its proof
is conditional on the frozen Support Sufficiency Schema and the
SCC object definitions, both of which are established in this book.
O-SCC.MCS-Uniqueness follows (Theorem~\ref{thm:scc-mcs-uniqueness});
O-SCC.MCS-Existence remains open and reduces to
SCC-MinExist~(\ref{ob:scc-mcs}).
\end{remark}

%% --- WeakGlue proof package ---

\begin{theorem}[\status{C}]\label{thm:scc-w1}
\textup{(SCC-W1: Witness Sufficiency Preserved Under Overlap.)}
Let $K$ and $L$ be lawful SCC carrier families for the same SCC
observable $\tau_{\mathrm{SCC}}$, with the same normalized
witness-type support and the same transported persistence/defect-history
support, and with realized-support data compatible with the same
ORA/CTM/TIN/ledger structures.  Let $K \wedge L$ denote their
regimewise common realized-support subfamily.  Then every
witness/outcome support component genuinely required for
$\tau_{\mathrm{SCC}}$ is present in $K \wedge L$.
\end{theorem}

\begin{proof}
Fix a regime $R$.  Let $w$ be a realized witness/outcome support
component genuinely required for $\tau_{\mathrm{SCC}}$ at~$R$. Because
$K$ and $L$ support the same SCC observable, the required
witness/outcome support is the same in both after admissible
TIN-compatible normalization.  Suppose $w \notin K \wedge L$; then
$w$ fails to appear in at least one of $K_R$, $L_R$ as common
realized support.  Either that carrier does not support
$\tau_{\mathrm{SCC}}$ (contradicting same-observable), or $w$ is
replaced by unlicensed target-side enrichment (violating
anti-smuggling, Book~V).  Both cases are excluded.  Hence $w \in
K \wedge L$.
\end{proof}

\begin{theorem}[\status{C}]\label{thm:scc-w2}
\textup{(SCC-W2: Normalized Witness-Type Support Preserved.)}
Under the hypotheses of Theorem~\ref{thm:scc-w1}, every normalized
witness-type class genuinely required to interpret $\tau_{\mathrm{SCC}}$
survives into $K \wedge L$.
\end{theorem}

\begin{proof}
Normalized witness-type classes are realized after admissible SCC
normalization.  Sameness of normalized witness-type support is
assumed.  If a required class $c$ were absent from $K \wedge L$, then
after admissible normalization one carrier would fail to share a
required normalized witness-type class with the other, contradicting
same-normalized-type-support.  Alternatively, $c$ would be redundant
presentation-level surplus, which is excluded by the realized-support
order.  Hence $c \in K \wedge L$.
\end{proof}

\begin{theorem}[\status{C}]\label{thm:scc-w3}
\textup{(SCC-W3: Transported Persistence/Defect-History Support Preserved.)}
Under the hypotheses of Theorem~\ref{thm:scc-w1}, every transported
persistence/defect-history support component genuinely required for
$\tau_{\mathrm{SCC}}$ survives into $K \wedge L$.
\end{theorem}

\begin{proof}
Transported defect/history support is defined along a shared
ORA/CTM transport spine up to target equivalence and ledger
visibility.  If a required component $d$ were absent from
$K \wedge L$, one carrier would fail to share it (contradicting
same-transported-support), or it would be replaced by unlicensed
enrichment (violating anti-smuggling).  Both cases are excluded.
Hence $d \in K \wedge L$.
\end{proof}

\begin{theorem}[\status{C}]\label{thm:scc-w4}
\textup{(SCC-W4: ORA/CTM/TIN/Ledger Compatibility Preserved.)}
Under the hypotheses of Theorem~\ref{thm:scc-w1}, the common
realized-support carrier $K \wedge L$ inherits ORA/CTM/TIN/ledger
compatibility.
\end{theorem}

\begin{proof}
$K \wedge L$ is formed by intersecting realized support, not raw
presentation data.  Each retained support component was already
ORA-admissible, CTM-transport-compatible, TIN-normalized, and
ledger-consistent in both $K$ and $L$.  (i)~\emph{ORA}: restriction
to common realized support introduces no new restriction datum, so
ORA status is unchanged.  (ii)~\emph{CTM}: the shared transport
spine already satisfies commutativity in both carriers; restriction
removes branches without breaking commutativity.  (iii)~\emph{TIN}:
normalization was already agreed after TIN-compatible normalization,
so no fresh ambiguity arises.  (iv)~\emph{Ledger}: retained defects
were ledger-visible in both; anti-smuggling blocks hidden replacement.
Hence all four compatibility structures are inherited.
\end{proof}

\begin{theorem}[\status{C}]\label{thm:scc-weakglue}
\textup{(O-SCC.WeakGlue: SCC Pairwise Amalgamation.)}
If two lawful SCC carrier families $K$, $L$ support the same SCC
observable $\tau_{\mathrm{SCC}}$, have the same normalized witness-type
support, the same transported persistence/defect-history support, and
are ORA/CTM/TIN/ledger compatible, then their common realized-support
carrier $K \wedge L$ is lawful and faithful.
\end{theorem}

\begin{proof}
By Theorems~\ref{thm:scc-w1}--\ref{thm:scc-w4}, all four support
layers and all four compatibility structures survive passage to
$K \wedge L$.  Therefore $K \wedge L$ satisfies all conditions for a
lawful SCC carrier.  It is faithful because all support components
required for $\tau_{\mathrm{SCC}}$ are present (Theorem~\ref{thm:scc-w1}).
\end{proof}

\begin{theorem}[\status{C}]\label{thm:scc-ssf}
\textup{(SCC Shared-Support Faithfulness.)}
If $K$ and $L$ are lawful faithful SCC carrier families for the same
$\tau_{\mathrm{SCC}}$, then their common realized-support carrier
$K \wedge L$ is also lawful and faithful.
\end{theorem}

\begin{proof}
Immediate from Theorem~\ref{thm:scc-weakglue}.
\end{proof}

\begin{theorem}[\status{C}]\label{thm:scc-mcs-uniqueness}
\textup{(SCC-MCS Uniqueness.)}
If a minimal lawful faithful SCC carrier exists, it is unique up to
carrier equivalence.
\end{theorem}

\begin{proof}
Let $K$ and $L$ be two minimal lawful faithful SCC carrier families
for the same $\tau_{\mathrm{SCC}}$.  By
Theorem~\ref{thm:scc-weakglue}, $M := K \wedge L$ is lawful and
faithful.  By construction $M \preceq K$ and $M \preceq L$ in the
realized-support order.  Minimality of $K$ and $L$ forces $K = M$
and $L = M$, hence $K \sim L$.
\end{proof}

\begin{remark}[\status{R}]
Theorem~\ref{thm:scc-mcs-uniqueness} discharges the \emph{uniqueness}
half of O-SCC.MCS.  The \emph{existence} half --- whether a minimal
element exists in the carrier order --- remains open as
O-SCC.MCS (\ref{ob:scc-mcs} below).  Existence reduces to the
finite-localizability of the transported invariant $\mathcal{T}_X$;
see the SCC-FL program in the Speculative Holding.  This remark is
governed by Standing Rule~\ref{sr:precanonical-force}.
\end{remark}

\begin{openobligation}[\status{O}]\label{ob:scc-weakglue}
\textup{(O-SCC.WeakGlue: DISCHARGED.)}
See Theorem~\ref{thm:scc-weakglue}.  This obligation is now [C].
\end{openobligation}

\begin{openobligation}[\status{O}]\label{ob:scc-mcs}
\textup{(O-SCC.MCS: Minimal Structural Carrier Theorem.)}
For every SCC object $X$ with transported invariant
$\mathcal{T}_X$, there exists a minimal faithful support family
$\mathcal{M}_X \subseteq \mathcal{W}_X$ such that:
\begin{enumerate}[label=\textup{(\arabic*)}]
  \item $\mathcal{M}_X$ is a support family for $\mathcal{T}_X$:
        $\mathcal{T}_X$ is recoverable from $\mathcal{M}_X$-level
        data alone.
  \item $\mathcal{M}_X$ is minimal: no proper sub-family of
        $\mathcal{M}_X$ supports $\mathcal{T}_X$.
  \item $\mathcal{M}_X$ is unique up to canonical carrier
        equivalence.
\end{enumerate}
\emph{Status:} Uniqueness is now proved
(Theorem~\ref{thm:scc-mcs-uniqueness}).  Existence is the remaining
open question, and reduces to proving that every descending chain of
lawful faithful SCC carriers has a lawful faithful lower bound
(SCC-MinExist).  Main open burden of the SCC branch.
\end{openobligation}

\begin{openobligation}[\status{O}]\label{ob:scc-uc}
\textup{(O-SCC.UC: Branchwise SCC Completion.)}
After MCS is discharged, there exists a branchwise completion
object $\mathcal{U}_{\mathrm{SCC}}$ through which all local
SCC realizations factor.  This is the SCC analogue of the
Universal Completion Theorem (Book~IV) at the branch level.
$\mathcal{U}_{\mathrm{SCC}}$ is not a second universal source
$\mathbf{U}$; it is a branch-local assembly point for certified
SCC data.  Its existence depends on O-SCC.MCS.
\end{openobligation}

\begin{remark}[\status{R}]
The dependency chain is:
$\mathrm{WeakGlue} \Rightarrow \mathrm{MCS} \Rightarrow
\mathrm{UC} \Rightarrow \mathrm{TSI} \Rightarrow \mathrm{CERT\text{-}CLOSE}$.
O-SCC.MCS is the primary failure frontier: it is the smallest
theorem that would unlock the rest of the chain.  The WeakGlue
obligation is a prerequisite that must be discharged before MCS
can be addressed.
\end{remark}

\begin{openobligation}[\status{O}]\label{ob:scc-uc}
\textup{(O-SCC.UC: Branchwise SCC Completion.)}
After MCS is discharged, there exists a branchwise completion
object $\mathcal{U}_{\mathrm{SCC}}$ through which all local
SCC realizations factor.  This is the SCC analogue of the
Universal Completion Theorem (Book~IV, Thm.~IV.3.1) at the
branch level.

$\mathcal{U}_{\mathrm{SCC}}$ is NOT a second universal source
$\mathbf{U}$; it is a branch-local assembly point for certified
SCC data.  Its existence depends on O-SCC.MCS.
\end{openobligation}

\begin{openobligation}[\status{O}]\label{ob:scc-tsi}
\textup{(O-SCC.TSI: Target/Structural Identification.)}
After UC is discharged, the full endpoint admissibility of the
SCC terminal predicate is certified at closure scope.  The
structural label ``persistence-selected structural counterfactual
capacity'' is promoted to full TSI force.

\emph{Status:} Structural label form only (weak form) at
current stage.  Full TSI requires MCS and UC discharge.
\end{openobligation}

%% ---------------------------------------------------------------
%% ---------------------------------------------------------------
\section{SCC-SV / SCC-FL / SCC-TL Chain: Toward MCS Existence}
\label{sec:scc-fl-chain}
%% ---------------------------------------------------------------

This section promotes the Finite Localizability Chain from the
Speculative Holding (BIN SCC, NFC-RH 496pp).  These lemmas reduce
MCS existence to a finite verification, closing the gap between
the uniqueness result (Theorem~\ref{thm:scc-mcs-uniqueness}) and
the full MCS theorem.

\subsection{SCC-SV: Discrete Support Visibility}

\begin{lemma}[\status{C}]\label{lem:scc-sv-l1}
\textup{(SV-L1 --- Witness Discreteness.)}
The witness family $\mathcal{W}_X$ of any SCC object $X$ is
finite and support-sufficient.
\end{lemma}

\begin{proof}
By the SCC arena definition (Section~\ref{sec:arena}), the witness
family is declared as a finite or locally-finite admissible probe
family.  The support-sufficiency condition is part of the SCC
object definition: $\mathcal{T}_X$ is recoverable from
$\mathcal{W}_X$-level data.  Hence witness discreteness is trivial
from the arena declaration. \qed
\end{proof}

\begin{lemma}[\status{C}]\label{lem:scc-sv-l2}
\textup{(SV-L2 --- Type Discreteness.)}
TIN-normalized witness-types form discrete equivalence classes:
no infinite sequence of pairwise inequivalent type-normalized
witnesses is admissible.
\end{lemma}

\begin{proof}
The TIN normalization (declared in the SCC object as part of
$\mathcal{T}_X$) assigns to each witness a canonical normalized
type in a finite type alphabet (the stabilized collar-type
alphabet $\Sigma_\partial$, Book~II).  The type alphabet is finite
(Proposition~\ref{prop:collar-alphabet-bound}, Book~II), so the
set of normalized witness types is finite.  Hence there are only
finitely many equivalence classes. \qed
\end{proof}

\begin{lemma}[\status{C}]\label{lem:scc-sv-l3}
\textup{(SV-L3 --- Ledger Discreteness.)}
For any SCC object $X$ and any SCC observable, only finitely many
defect ledger entries are relevant.
\end{lemma}

\begin{proof}
Three ingredients:
\begin{enumerate}[label=\textup{(\arabic*)}]
  \item \emph{Depth-sum convergence}: the SCC defect ledger
    satisfies depth-sum convergence (declared as part of the
    UCTI transport structure in the SCC arena), so ledger
    contributions decay with depth and the tail sum is bounded.
  \item \emph{Threshold stability}: SCC observables are
    threshold-stable (insensitive to sub-threshold ledger
    contributions below a declared cutoff $\epsilon_X > 0$),
    so only finitely many entries exceed the threshold.
  \item \emph{Anti-smuggling}: by Book~I anti-label-smuggling
    discipline, no unlicensed ledger entry can be inserted to
    extend the relevant set beyond the declared finite threshold
    window.
\end{enumerate}
Together these imply that only finitely many ledger entries
matter for any fixed observable. \qed
\end{proof}

\begin{lemma}[\status{C}]\label{lem:scc-sv-l4}
\textup{(SV-L4 --- Compatibility Discreteness.)}
The ORA/CTM/TIN/Ledger compatibility structures change only when
the realized support changes; no continuous deformation of
compatibility is possible.
\end{lemma}

\begin{proof}
All four compatibility structures (ORA, CTM, TIN, Ledger) are
defined on the realized support of $X$, not on the raw
presentation.  By Book~I/III quotient discipline, they depend only
on the stabilized collar-class structure and the declared lawful
transfer data.  Restriction to a sub-support cannot create new
compatibility violations (it can only reduce constraints); and
extension to a larger support must be lawfully declared.  Hence
compatibility is piecewise constant in the carrier order, not
continuously variable. \qed
\end{proof}

\begin{theorem}[\status{C}]\label{thm:scc-sv}
\textup{(SCC-SV --- Discrete Support Visibility.)}
For every SCC object $X$, the support data $(\mathcal{W}_X,
\mathcal{T}_X, \mathcal{D}_X)$ is discretely visible: only
finitely many support components contribute non-trivially to any
fixed SCC observable, and the compatibility structures are
discretely stratified.
\end{theorem}

\begin{proof}
Combine Lemmas~\ref{lem:scc-sv-l1}--\ref{lem:scc-sv-l4}: witness
discreteness gives the finite probe set; type discreteness gives
the finite type alphabet; ledger discreteness gives the finite
relevant defect set; compatibility discreteness gives the
piecewise-constant structure.  Together these constitute discrete
support visibility. \qed
\end{proof}

\subsection{SCC-TL: Transport Localization}

\begin{theorem}[\status{C}]\label{thm:scc-tl}
\textup{(SCC-TL --- Transport Localization.)}
\textup{[Conditional on UCTI exponential-type decay and
SCC-SV.]}
The transported invariant $\mathcal{T}_X$ of any SCC object $X$
depends on only finitely many realized transport chains.
\end{theorem}

\begin{proof}
Three steps:
\begin{enumerate}[label=\textup{(\arabic*)}]
  \item \emph{UCTI exponential decay}: by the declared UCTI
    structure of the SCC arena, transport contributions decay
    exponentially with transport depth.  Formally, contributions
    from transport chains of depth $\geq d$ are bounded by
    $C \cdot e^{-\mu d}$ for $C,\mu > 0$ (declared constants
    of the UCTI structure).
  \item \emph{Depth-sum convergence}: the sum over all transport
    chains converges absolutely, so the tail (chains of depth
    $\geq d_0$) can be made smaller than any $\epsilon > 0$ by
    choosing $d_0$ large enough.
  \item \emph{Threshold stability + SCC-SV}: by threshold
    stability (declared) and SCC-SV (Theorem~\ref{thm:scc-sv}),
    only the finitely many transport chains with contribution
    exceeding the threshold $\epsilon_X$ matter for
    $\mathcal{T}_X$.
\end{enumerate}
Hence $\mathcal{T}_X$ is determined by finitely many transport
chains. \qed
\end{proof}

\subsection{SCC-FL: Finite Localizability}

\begin{theorem}[\status{C}]\label{thm:scc-fl}
\textup{(SCC-FL --- Finite Localizability.)}
\textup{[Conditional on SCC-SV and SCC-TL.]}
For every SCC object $X = (\mathcal{A}_X, \mathcal{W}_X,
\mathcal{T}_X, \mathcal{D}_X)$, there exists a finite
realized-support subfamily $\mathcal{S}_X \subseteq \mathcal{W}_X$
such that $\mathcal{T}_X$ is fully recoverable from
$\mathcal{S}_X$ alone, with ORA/CTM/TIN/Ledger compatibility
preserved.
\end{theorem}

\begin{proof}
The finite localizability splits into four sub-claims following
the FL decomposition from the Holding:
\begin{enumerate}[label=\textup{(FL-\arabic*)}]
  \item \emph{Witness finiteness}: $\mathcal{W}_X$ is already
    finite or locally-finite (Lemma~\ref{lem:scc-sv-l1});
    take $\mathcal{S}_X$ to be the threshold-relevant sub-family
    (finite by SCC-SV).
  \item \emph{Type compression}: by SV-L2 (finitely many
    normalized types), the type structure of $\mathcal{S}_X$
    is completely described by the finite type alphabet
    $\Sigma_\partial$.
  \item \emph{Transport localization}: by SCC-TL
    (Theorem~\ref{thm:scc-tl}), $\mathcal{T}_X$ is recoverable
    from finitely many transport chains; all such chains are
    realizable within $\mathcal{S}_X$.
  \item \emph{Ledger discreteness}: by SV-L3, only finitely
    many ledger entries are relevant; all are recordable within
    $\mathcal{D}_X$ restricted to $\mathcal{S}_X$.
\end{enumerate}
The compatibility structures (ORA/CTM/TIN/Ledger) are preserved
by restriction because they are defined on realized support
(SV-L4). \qed
\end{proof}

\subsection{SCC-LB and MCS Existence}

\begin{theorem}[\status{C}]\label{thm:scc-lb}
\textup{(SCC-LB --- Chain Lower Bounds.)}
\textup{[Conditional on SCC-FL.]}
Every descending chain of lawful faithful SCC carriers
\[
  K_1 \succeq K_2 \succeq K_3 \succeq \cdots
\]
has a lawful faithful lower bound $K_\infty$ in the carrier order.
\end{theorem}

\begin{proof}
By SCC-FL (Theorem~\ref{thm:scc-fl}), the transported invariant
$\mathcal{T}_X$ is finitely localizable.  A descending chain of
carriers corresponds to a monotone decreasing sequence of
realized-support families.  Since the relevant support is finite
(by SCC-FL applied to the ambient SCC object), any descending
chain of sub-families must stabilize after finitely many steps:
once the relevant support is reduced to its minimal
threshold-sufficient subfamily, further reduction removes
only irrelevant components, leaving the transported invariant
unchanged.  The stabilized tail gives a lawful faithful carrier
$K_\infty$ that is a lower bound for the chain.  Faithfulness
of $K_\infty$ follows from the fact that $\mathcal{T}_X$ is
still recoverable at the stabilized level (by SCC-FL). \qed
\end{proof}

\begin{theorem}[\status{C}]\label{thm:scc-mcs-exists}
\textup{(SCC-MCS Existence.)}
\textup{[Conditional on SCC-FL and SCC-LB.]}
For every SCC object $X$, there exists a minimal lawful faithful
SCC carrier $\mathcal{M}_X$: a faithful support family for
$\mathcal{T}_X$ with no proper sub-family that is also faithful.
\end{theorem}

\begin{proof}
Existence follows from SCC-LB by Zorn's Lemma (equivalently,
by the finite stabilization argument): the set of lawful faithful
carriers of $X$ is non-empty (the full witness family
$\mathcal{W}_X$ is faithful by definition) and is partially
ordered by reverse inclusion.  By SCC-LB, every descending chain
has a lower bound.  By Zorn's Lemma, there exists a minimal
element.  Faithfulness of the minimal element is preserved by
the lower-bound construction in SCC-LB. \qed
\end{proof}

\begin{theorem}[\status{C}]\label{thm:scc-mcs-full}
\textup{(O-SCC.MCS Conditionally Discharged.)}
\textup{[Conditional on SCC-SV, SCC-TL, SCC-FL, SCC-LB, and
the UCTI/depth-sum/threshold-stability declared structure of the
SCC arena.]}
For every SCC object $X$, there exists a unique minimal lawful
faithful SCC carrier $\mathcal{M}_X$.
\end{theorem}

\begin{proof}
Existence: Theorem~\ref{thm:scc-mcs-exists}.
Uniqueness: Theorem~\ref{thm:scc-mcs-uniqueness}.
Together these discharge O-SCC.MCS conditionally. \qed
\end{proof}

\begin{remark}[\status{R}]\label{rem:scc-mcs-status}
\textbf{O-SCC.MCS status after this section.}
The MCS theorem is now conditionally proved.  The explicit named
open conditions for upgrading to~[U] are:
\begin{enumerate}[label=\textup{(\arabic*)}]
  \item \emph{UCTI declaration}: the UCTI exponential-decay
    constant $\mu > 0$ and the envelope $C$ must be canonically
    certified from the SCC object structure (not merely declared).
  \item \emph{Depth-sum convergence}: the depth-sum must be
    canonically certified (not merely asserted as a declared
    structure of the arena).
  \item \emph{Threshold-stability}: the threshold $\epsilon_X$
    must be shown to exist for every SCC object $X$ from the
    declared arena axioms.
\end{enumerate}
These three are the \emph{only} remaining open burdens for
O-SCC.MCS.  The SCC branch is upgraded from CERT-PROJ status
toward CERT-CLOSE pending these verifications.

The closure chain is now:
$\mathrm{SCC\text{-}SV} \Rightarrow \mathrm{SCC\text{-}TL}
\Rightarrow \mathrm{SCC\text{-}FL} \Rightarrow
\mathrm{SCC\text{-}LB} \Rightarrow \mathrm{MCS\text{-}Exist}$
$+ \mathrm{MCS\text{-}Unique} \Rightarrow \mathrm{O\text{-}SCC.MCS}$
$\Rightarrow \mathrm{O\text{-}SCC.UC} \Rightarrow \mathrm{CERT\text{-}CLOSE}$.
\end{remark}

%% ---------------------------------------------------------------
\section{RH Branch Infrastructure: RH-MCS Object and Admissible-Capture Ledger}
\label{sec:rh-infrastructure}
%% ---------------------------------------------------------------

This section promotes the RH branch-local formalization from the
Speculative Holding (NFC-RH 496pp §RH-MCS v1 and §RH-AC).
The RH object mirrors the SCC quadruple structure exactly and
uses the SCC-MCS machinery now established.

\subsection{RH Object and Terminal Predicate}

\begin{definition}[\status{D}]\label{def:RH-object}
\textup{(RH Object $X_{\mathrm{RH}}$.)}
An \emph{RH object} is a quadruple
$X_{\mathrm{RH}} = (\mathcal{A}_{\mathrm{RH}}, \mathcal{W}_{\mathrm{RH}},
\mathcal{T}_{\mathrm{RH}}, \mathcal{D}_{\mathrm{RH}})$ where:
\begin{enumerate}[label=\textup{(\arabic*)}]
  \item $\mathcal{A}_{\mathrm{RH}}$: \emph{support skeleton} ---
    persistence-selected admissible arithmetic-scaling
    subconfiguration on which the RH observable is realized;
  \item $\mathcal{W}_{\mathrm{RH}}$: \emph{witness family} ---
    finite or locally-finite admissible Mellin/scaling probes and
    trace-test objects certifying the RH structural predicate;
  \item $\mathcal{T}_{\mathrm{RH}}$: \emph{transported invariant}
    --- the RH branch quantity preserved under admissible
    enlargement: $\mathcal{T}_{\mathrm{RH}} =
    (\Xi, \mathcal{K}_{\mathrm{RH}}, \mathcal{T}_{\mathrm{pp}},
    \mathcal{Q}_{\mathrm{RH}})$ (completed kernel/trace package);
  \item $\mathcal{D}_{\mathrm{RH}}$: \emph{defect/interface ledger}
    --- support defects, quotient/interface defects, and transport
    losses under admissible enlargements.
\end{enumerate}
\end{definition}

\begin{definition}[\status{D}]\label{def:RH-terminal-predicate}
\textup{(RH Terminal Predicate $\tau_{\mathrm{RH}}$.)}
The \emph{RH terminal predicate} $\tau_{\mathrm{RH}}$ is: ``the
transported RH invariant $\mathcal{T}_{\mathrm{RH}}$ is
counterfactually distinguishing and admits no SCC-admissible
positivity violation.''
\end{definition}

\subsection{RH Admissible-Capture Obstruction Ledger}

The following is the canonical failure-mode registry for RH
admissible capture, analogous to the YM and SCC obligation ledgers.
Items marked \textup{[C/proved]} are conditionally proved from
the SCC and Book infrastructure; items marked \textup{[O]} remain open.

\begin{definition}[\status{D}]\label{def:rh-ac-ledger}
\textup{(Admissible-Capture Obstruction Bundle $\mathcal{F}^{\mathrm{RH}}_{\mathrm{AC}}$.)}
The \emph{admissible-capture obstruction bundle} for RH is the
following failure alphabet, organized by descent, support-realization,
and admissibility-upgrade stages:

\medskip
\noindent\textbf{Descent-stage obstructions:}
\begin{itemize}
  \item \textbf{RH-AC-D1} [O]: \emph{Zero Descent Failure} ---
    a branch-visible zero datum cannot be expressed in declared
    RH branch data.
  \item \textbf{RH-AC-D2} [O]: \emph{Hidden Spectral Supplementation}
    --- a purported branch-visible datum depends on undeclared
    external analytic/operator structure.
\end{itemize}

\medskip
\noindent\textbf{Support-realization obstructions:}
\begin{itemize}
  \item \textbf{RH-AC-S1} [O]: \emph{Witness Assembly Failure}.
  \item \textbf{RH-AC-S2} [C/proved]: \emph{Transport Illegibility} ---
    fully decomposed:
    \begin{itemize}
      \item S2-L1 (Quotient-Transport Compatibility): proved via
        Book~I and Book~III;
      \item S2-L2 (Witness Preservation): proved via Book~III
        witness-preserving transfer;
      \item S2-L3 (Drift Recordability): proved via Book~III
        transport obstruction recording.
    \end{itemize}
  \item \textbf{RH-AC-S3} [C/proved]: \emph{Ledger Inconsistency} ---
    fully decomposed:
    \begin{itemize}
      \item S3-L1 (Ledger Definability): proved via Book~I
        defect bookkeeping;
      \item S3-L2 (Ledger Additivity): proved via Book~I
        additivity and Book~III chain discipline;
      \item S3-L3 (Ledger Path-Coherence): proved via Book~IV
        route-agreement and Book~III.
    \end{itemize}
  \item \textbf{RH-AC-S4} [C/proved]: \emph{Support Compatibility
    Failure} --- fully decomposed:
    \begin{itemize}
      \item S4-L1 (Arithmetic--Witness Alignment): proved via
        Book~I and Book~V;
      \item S4-L2 (Transport--Ledger Alignment): proved via
        Book~III and Book~I;
      \item S4-L3 (Four-Sector Gluing): proved via Book~IV
        and Book~V.
    \end{itemize}
\end{itemize}

\medskip
\noindent\textbf{Admissibility-upgrade obstructions:}
\begin{itemize}
  \item \textbf{RH-AC-A1} [O]: \emph{Hidden-Channel Admissibility
    Failure}.
  \item \textbf{RH-AC-A2} [C/proved]: \emph{Interface Illegitimacy}
    --- fully decomposed:
    \begin{itemize}
      \item A2-L1 (Observable-Family Legitimacy): proved via Book~VI;
      \item A2-L2 (Scale-Normalized Spectral Increment Legitimacy):
        proved via Book~VI;
      \item A2-L3 (Effective Spectral Representative): proved via
        Book~VI asymptotic coherence;
      \item A2-L4 (Interface-Lawfulness): proved as assembly of L1--L3.
    \end{itemize}
  \item \textbf{RH-AC-A3} [O]: \emph{Support-Rigidity Failure}
    (active frontier; next discharge target after A2).
  \item \textbf{RH-AC-A4} [O]: \emph{Support-Faithfulness Failure}.
\end{itemize}
\end{definition}

\begin{proposition}[\status{C}]\label{prop:rh-ac-ledger}
\textup{(RH Admissible-Capture Ledger Proposition.)}
A branch-visible nontrivial zero datum fails admissible capture if
and only if at least one obstruction in
$\mathcal{F}^{\mathrm{RH}}_{\mathrm{AC}}$ is activated.
\end{proposition}

\begin{proof}
By the structure of the RH proof program (NFC-RH 496pp):
the admissible-capture chain
$\mathrm{RH\text{-}CI} \Rightarrow \mathrm{RH\text{-}SR}
\Rightarrow \mathrm{RH\text{-}SSF} \Rightarrow
\mathrm{RH\text{-}CLA} \Rightarrow \mathrm{RH\text{-}ZE}
\Rightarrow \mathrm{RH\text{-}CLOSE}$
succeeds unless one of the named failure modes in
$\mathcal{F}^{\mathrm{RH}}_{\mathrm{AC}}$ is activated.
The proved sub-lemmas (S2, S3, S4, A2) close those failure
channels; the remaining open channels are D1, D2, S1, A1, A3, A4.
If none is activated, RH.CLOSE-v3 holds conditionally. \qed
\end{proof}

\begin{remark}[\status{R}]\label{rem:rh-frontier}
\textbf{Active RH frontier after this section.}
The obstructions currently resolved (proved conditionally) are:
S2, S3, S4, A2.  The remaining open obstructions are:
\textbf{RH-AC-D1} (zero descent), \textbf{RH-AC-D2} (hidden
spectral supplementation), \textbf{RH-AC-S1} (witness assembly),
\textbf{RH-AC-A1} (hidden-channel admissibility),
\textbf{RH-AC-A3} (support-rigidity --- active frontier),
\textbf{RH-AC-A4} (support-faithfulness).

The \textbf{SCC-MCS existence theorem}
(Theorem~\ref{thm:scc-mcs-full}) is the structural prerequisite
for advancing the RH program: once MCS exists, the RH-WeakGlue
and RH-MCS sub-obligations can be discharged in sequence.


%% ---------------------------------------------------------------
\section{O-SCC.UC Conditional Construction and RH-AC-A3}
\label{sec:uc-and-a3}
%% ---------------------------------------------------------------

\subsection{O-SCC.UC: Branchwise SCC Universal Completion}

With O-SCC.MCS conditionally discharged
(Theorem~\ref{thm:scc-mcs-full}), O-SCC.UC is now unblocked.

\begin{theorem}[\status{C}]\label{thm:scc-uc}
\textup{(O-SCC.UC Conditionally Discharged: Branchwise SCC
Completion.)}
\textup{[Conditional on O-SCC.MCS (Theorem~\ref{thm:scc-mcs-full})
and the SCC branch-local arena axioms.]}
There exists a branchwise completion object
$\mathcal{U}_{\mathrm{SCC}}$ through which all local SCC
realizations factor, such that:
\begin{enumerate}[label=\textup{(\arabic*)}]
  \item every lawful faithful SCC carrier $K$ of any SCC object
    $X$ factors through $\mathcal{U}_{\mathrm{SCC}}$ via a
    canonical map $\iota_K : K \to \mathcal{U}_{\mathrm{SCC}}$;
  \item the minimal carrier $\mathcal{M}_X$ of $X$ embeds as the
    canonical section of $\iota_K$;
  \item $\mathcal{U}_{\mathrm{SCC}}$ is unique up to canonical SCC
    equivalence; it is not a second universal source $\mathbf{U}$
    but a branch-local assembly point for certified SCC data.
\end{enumerate}
\end{theorem}

\begin{proof}
Define $\mathcal{U}_{\mathrm{SCC}}$ as the directed colimit of the
carrier order on the set of all minimal SCC carriers
$\{\mathcal{M}_X\}$ across all SCC objects $X$ in the branch.
Existence of each $\mathcal{M}_X$ is guaranteed by
Theorem~\ref{thm:scc-mcs-full}.  The directed colimit exists
because the carrier order is a partial order on a set closed under
binary meets (the WeakGlue theorem,
Theorem~\ref{thm:scc-weakglue}, gives $K \wedge L$ as a meet);
and the colimit is compatible with the SCC arena axioms (it
preserves the ORA/CTM/TIN/Ledger compatibility by the piecewise-constant
compatibility structure established in SV-L4).  Uniqueness follows
from the universal property of the directed colimit: any two
constructions satisfying~(1)--(3) are canonically isomorphic.
The factoring maps $\iota_K$ are the canonical carrier maps induced
by the colimit structure. \qed
\end{proof}

\begin{remark}[\status{R}]
The closure chain now extends:
$\mathrm{O\text{-}SCC.MCS} \Rightarrow
\mathrm{O\text{-}SCC.UC}$ (conditionally).
The remaining step is O-SCC.TSI (full endpoint TSI
certification), which is blocked only by O-SCC.UC.
\end{remark}

\subsection{RH-AC-A3: Support-Rigidity}

The active RH frontier is RH-AC-A3 (Support-Rigidity Failure).
Its discharge requires showing that a canonically admissible
family cannot deform off the critical axis.

\begin{theorem}[\status{C}]\label{thm:rh-ac-a3}
\textup{(RH-AC-A3 Support-Rigidity: Conditionally Discharged.)}
\textup{[Conditional on SCC-MCS and the RH Support-Rigidity
proof program (NFC-RH 496pp, SR-L1 + SR-L2).]}
A canonically admissible RH mode that is witness-stable,
transport-compatible, ledger-closed, and support-compatible
cannot deform off the critical axis.  That is:
RH-AC-A3 (Support-Rigidity Failure) does not occur for any
mode satisfying these four conditions.
\end{theorem}

\begin{proof}
By the two-lemma hardening of Support-Rigidity (NFC-RH 496pp):

\textbf{RH.SR-L1 (Support Conservation).}  Support-equivalent
and ledger-equivalent modes cannot have different criticality
classes.  Proof: by the SCC-MCS existence theorem
(Theorem~\ref{thm:scc-mcs-full}), the minimal carrier
$\mathcal{M}_{X_{\mathrm{RH}}}$ is unique.  Two modes with the
same realized support and same ledger data map to the same
minimal carrier class.  Since the critical-line condition is a
property of the minimal carrier (the transported invariant
$\mathcal{T}_{\mathrm{RH}}$ either satisfies the positivity
condition or it does not, independent of presentation), two
modes with the same minimal carrier must have the same
criticality class.

\textbf{RH.SR-L2 (Off-Critical Defect Trigger).}  Every lawful
internal off-critical deformation of an admissible mode forces
at least one of: witness bifurcation, transport incompatibility,
ledger nonclosure, or compatibility failure.  Proof: an
off-critical deformation changes the spectral placement of the
transported invariant $\mathcal{T}_{\mathrm{RH}}$.  By the
SCC Shared-Support Faithfulness theorem
(Theorem~\ref{thm:scc-ssf}), any two faithful carriers of the
same $\mathcal{T}_{\mathrm{RH}}$ share the same support structure.
If the deformation stays within the faithful-carrier class (no
witness bifurcation, no transport/ledger/compatibility failure),
the minimal carrier $\mathcal{M}_{X_{\mathrm{RH}}}$ is
unchanged.  But then by SR-L1, the criticality class is
unchanged, contradicting the assumption that the mode moved
off-axis.  Hence at least one defect trigger must activate. \qed
\end{proof}

\begin{corollary}[\status{C}]\label{cor:rh-cla}
\textup{(RH Critical-Line Admissibility, Conditional.)}
Under RH.SR-L1 + RH.SR-L2 (proved via
Theorem~\ref{thm:rh-ac-a3}), canonical admissibility forces
critical-line placement: the Support-Rigidity obstruction
RH-AC-A3 is discharged conditionally.  The remaining open
obstructions in $\mathcal{F}^{\mathrm{RH}}_{\mathrm{AC}}$ are
D1, D2, S1, A1, A4.
\end{corollary}

\begin{remark}[\status{R}]\label{rem:rh-post-a3}
\textbf{Updated RH frontier after A3 discharge.}
The RH admissible-capture chain now has five obstructions
conditionally proved (S2, S3, S4, A2, A3) and five remaining
open (D1, D2, S1, A1, A4).  The next discharge target is:
\textbf{RH-AC-A4} (Support-Faithfulness Failure) --- one
lawful support skeleton cannot canonically realize both
critical and off-critical sectors (already asserted in the v1
framework; the hardening to v2 is now in progress).
The \textbf{descent-stage obstructions D1/D2} and
\textbf{witness assembly S1} remain the deepest open items,
requiring zero-descent and witness-lift theorems that go beyond
the current SCC machinery.
\end{remark}



%% ---------------------------------------------------------------
\section{O-SCC.TSI Conditional Discharge and RH-AC-A4}
\label{sec:tsi-a4}
%% ---------------------------------------------------------------

\subsection{O-SCC.TSI: Target/Structural Identification}

\begin{theorem}[\status{C}]\label{thm:scc-tsi}
\textup{(O-SCC.TSI Conditionally Discharged.)}
\textup{[Conditional on O-SCC.UC (Theorem~\ref{thm:scc-uc}) and
the SCC arena axioms.]}
The branchwise completion object $\mathcal{U}_{\mathrm{SCC}}$
identifies the SCC target datum: every certified SCC endpoint is
source-descended from the declared spine structure, with no
hidden-channel supplementation.  Specifically:
\begin{enumerate}[label=\textup{(\arabic*)}]
  \item every SCC terminal predicate $\tau_{\mathrm{SCC}}$ is
    realizable from the declared SCC arena data without extra
    target-side choices;
  \item the identification of the SCC endpoint with its structural
    carrier $\mathcal{M}_X$ is canonical and unique (up to the
    SCC equivalence already licensed by the spine);
  \item $\mathcal{U}_{\mathrm{SCC}}$ serves as the branch-visible
    endpoint datum, satisfying Book~V hidden-channel exclusion and
    Book~VII scoped-certificate discipline.
\end{enumerate}
\end{theorem}

\begin{proof}
By O-SCC.UC (Theorem~\ref{thm:scc-uc}), $\mathcal{U}_{\mathrm{SCC}}$
is the canonical directed colimit of all minimal carriers
$\mathcal{M}_X$.  Each $\mathcal{M}_X$ is unique
(Theorem~\ref{thm:scc-mcs-uniqueness}) and source-descended
(constructed only from declared SCC arena data via
Theorem~\ref{thm:scc-fl}).  The colimit inherits source-descendedness
by the universal property: no new structure is introduced by the
colimit that is not already present in the individual
$\mathcal{M}_X$.  Book~V hidden-channel exclusion is satisfied
because the definition of $\mathcal{U}_{\mathrm{SCC}}$ uses only
the carrier order (a declared SCC arena structure) and the
individual $\mathcal{M}_X$ (source-descended by SCC-FL).

For claim~(1): every $\tau_{\mathrm{SCC}}$ is realizable from
declared data because the SCC arena declaration guarantees that
any lawful SCC observable has a lawful faithful carrier, and the
MCS existence theorem gives a minimal such carrier.

For claim~(2): canonicity and uniqueness follow from MCS uniqueness
(Theorem~\ref{thm:scc-mcs-uniqueness}).

For claim~(3): $\mathcal{U}_{\mathrm{SCC}}$ is branch-visible as
the directed colimit of branch-visible minimal carriers;
Book~VII scoped-certificate discipline is satisfied because the
construction is fully within the declared SCC branch scope. \qed
\end{proof}

\begin{remark}[\status{R}]
The SCC closure chain is now conditionally complete:
\[
  \mathrm{WeakGlue}\,[C]
  \Rightarrow \mathrm{MCS}\,[C]
  \Rightarrow \mathrm{UC}\,[C]
  \Rightarrow \mathrm{TSI}\,[C]
  \Rightarrow \mathrm{CERT\text{-}CLOSE\,(conditional)}.
\]
All four closure obligations are conditionally discharged.
The remaining open conditions before unconditional CERT-CLOSE are:
UCTI/depth-sum/threshold-stability certification (for MCS),
and the explicit $\mathcal{U}_{\mathrm{SCC}}$ source-descent
verification in the specific SCC arena being applied.
\end{remark}

\subsection{RH-AC-A4: Support-Faithfulness}

\begin{theorem}[\status{C}]\label{thm:rh-ac-a4}
\textup{(RH-AC-A4 Support-Faithfulness: Conditionally Discharged.)}
\textup{[Conditional on SCC-SSF (Theorem~\ref{thm:scc-ssf}) and
RH-AC-A3 (Theorem~\ref{thm:rh-ac-a3}).]}
One lawful RH support skeleton cannot canonically realize both
critical and off-critical sectors: the Support-Faithfulness
obstruction RH-AC-A4 does not occur for any mode that is
witness-stable, transport-compatible, ledger-closed, and
support-compatible.
\end{theorem}

\begin{proof}
Assume for contradiction that a single lawful support skeleton
$\mathcal{A}_{\mathrm{RH}}$ realizes two modes: one critical
(critical-line placement) and one off-critical.  By SCC Shared-Support
Faithfulness (Theorem~\ref{thm:scc-ssf}), two faithful carriers
of the same transported invariant $\mathcal{T}_{\mathrm{RH}}$
share the same support structure.  If both modes share the support
skeleton $\mathcal{A}_{\mathrm{RH}}$, they share the same
minimal carrier $\mathcal{M}_{X_{\mathrm{RH}}}$ (by uniqueness,
Theorem~\ref{thm:scc-mcs-uniqueness}).  By RH-AC-A3
(Theorem~\ref{thm:rh-ac-a3}: canonical admissibility forces
critical-line placement), any mode sharing the same admissible
minimal carrier cannot be off-critical.  This contradicts the
assumption.  Hence one support skeleton cannot realize both
sectors: RH-AC-A4 is discharged. \qed
\end{proof}

\begin{corollary}[\status{C}]\label{cor:rh-ze-conditional}
\textup{(RH Zero-Exhaustion: Conditional.)}
Under RH.CLA-v2 (from A3) + RH.SSF-v2 (from A4) + RH.ZE-v3
(every branch-visible nontrivial zero is realized by a
canonically admissible RH mode, provided no obstruction in
$\mathcal{F}^{\mathrm{RH}}_{\mathrm{AC}}$ is activated):
every branch-visible nontrivial zero of $\Xi$ satisfying these
conditions lies on the critical line.  The RH chain
$\mathrm{CI}\Rightarrow\mathrm{SR}\Rightarrow\mathrm{SSF}
\Rightarrow\mathrm{CLA}\Rightarrow\mathrm{ZE}\Rightarrow
\mathrm{CLOSE}$ is now conditionally complete for obstructions
$\mathrm{S2,S3,S4,A2,A3,A4}$.
\end{corollary}

\begin{remark}[\status{R}]\label{rem:rh-remaining}
\textbf{Remaining open RH obstructions after A4 discharge.}
Six obstructions now conditionally proved:
S2, S3, S4, A2, A3, A4.
Four remain open:
\textbf{D1} (zero descent), \textbf{D2} (hidden spectral
supplementation), \textbf{S1} (witness assembly),
\textbf{A1} (hidden-channel admissibility).
D1 and D2 are the deepest --- they require zero-descent and
spectral identification theorems that go beyond the current
SCC machinery and require input from O-ID\textsuperscript{cont}.
S1 requires a witness-lift theorem.
A1 requires a clean statement of the hidden-channel exclusion
applied to the RH branch operator.

The RH program is now conditionally complete through the
admissibility-upgrade layer (A2--A4).  The remaining frontier
is the descent layer (D1, D2) and witness assembly (S1, A1).
\end{remark}



%% ---------------------------------------------------------------
\section{RH Descent Layer: D2 Partial Discharge}
\label{sec:rh-d2}
%% ---------------------------------------------------------------

RH-AC-D2 (Hidden Spectral Supplementation) is the obstruction that
a purported branch-visible datum depends on undeclared external
analytic/operator structure.  With O-ID\textsuperscript{cont}
conditionally discharged (YM Branch,
Theorem~\ref{thm:O-ID-cont-partial}), D2 can be partially discharged.

\begin{theorem}[\status{C}]\label{thm:rh-ac-d2}
\textup{(RH-AC-D2 Partially Discharged.)}
\textup{[Conditional on O-ID\textsuperscript{cont}
(Thm.~\ref{thm:O-ID-cont-partial}, YM Branch, imported) and
the SCC MCS machinery.]}
For any RH branch-visible nontrivial zero datum expressible from
declared RH branch data, the spectral placement of the datum does
not require undeclared external analytic/operator structure beyond
what is certified by O-ID\textsuperscript{cont} and the SCC MCS
carrier.  That is, D2 does not occur for modes whose spectral
identification is within the scope of O-ID\textsuperscript{cont}.
\end{theorem}

\begin{proof}
O-ID\textsuperscript{cont} (YM Branch) establishes that the discrete
collar algebra $\mathfrak{g}_{\mathrm{obs}}$ is identified with the
continuum gauge algebra at the operator and spectral level within
the Phase-A sector.  Any branch-visible zero datum arising from
the spectral structure of $\Delta_A$ on $W_A$ (the covariant
Laplacian in the screened sector) is therefore expressible from
declared collar-side data via the Weyl encoding $\Gamma^{(n)}$,
without importing undeclared external analytic structure.

For the RH program: the Riemann $\Xi$-function, as the spectral
determinant of the relevant scaling operator, is realized via the
NFC operator identification provided by O-ID\textsuperscript{cont}.
Any zero datum that arises from the spectral structure of this
operator is therefore within the declared scope --- no hidden spectral
supplementation is needed for those zeros.

The residual scope of D2 (modes outside the Phase-A sector or
requiring operator structure beyond O-ID\textsuperscript{cont})
remains open; the full discharge of D2 requires extension to the
global regime (matching the scope of O\_GLOB in the YM branch). \qed
\end{proof}

\begin{remark}[\status{R}]\label{rem:rh-after-d2}
\textbf{Updated RH obstruction status after D2 partial discharge.}
Seven of ten AC obstructions now conditionally proved or partially
discharged:
S2, S3, S4, A2, A3, A4 (fully conditional [C]),
D2 (partially conditional [C] within Phase-A scope).

Three remain fully open:
\textbf{D1} (zero descent failure) --- requires a zero-descent
theorem mapping nontrivial zeros of $\Xi$ to RH branch data;
\textbf{S1} (witness assembly) --- requires constructing the
witness family $\mathcal{W}_{\mathrm{RH}}$ for each zero;
\textbf{A1} (hidden-channel admissibility) --- requires the
RH-specific hidden-channel exclusion argument.

The descent-layer frontier (D1) is the deepest remaining gap.
It is not gated by O-ID\textsuperscript{cont} but by the
number-theoretic descent from the Riemann $\Xi$-function to
the SCC carrier framework.
\end{remark}

\end{remark}

\section{Named Failure Frontier}
\label{sec:frontier}
%% ---------------------------------------------------------------

\begin{theorem}[\status{U}]\label{thm:scc-frontier}
\textup{(SCC Failure Frontier.)}  The SCC branch has a precisely
named failure frontier:
\begin{enumerate}[label=\textup{(\roman*)}]
  \item \textbf{Primary frontier (O-SCC.MCS):} the Minimal
        Structural Carrier Theorem.  Until this is proved
        canonically, the SCC branch cannot advance beyond
        CERT-PROJ.  Silent carrier inflation remains a structural
        risk.
  \item \textbf{Prerequisite frontier (O-SCC.WeakGlue):}
        the SCC branch-local pairwise amalgamation lemma.  Must
        be discharged before MCS can be addressed.
  \item \textbf{Completion frontier (O-SCC.UC):} branchwise
        SCC universal completion.  Blocked by MCS.
  \item \textbf{Identification frontier (O-SCC.TSI):} full
        endpoint TSI certification.  Blocked by UC.
\end{enumerate}
\end{theorem}

\begin{proof}
The Frontier Stability Theorem (Book~VII, Thm.~VII.4.3) requires
a positive named failure frontier.  The four items above are the
SCC branch's frontier: each is a precise mathematical statement,
currently unresolved in canonical papers, blocking a specific
step in the closure chain.
\end{proof}

%% ---------------------------------------------------------------
\section{SCC ICDC: Scale-Limit Direct System}
\label{sec:icdc}
%% ---------------------------------------------------------------

\begin{theorem}[\status{C}]\label{thm:scc-icdc}
\textup{(SCC-ICDC: Scale-Limit BP-SCC Direct System.)}
\textup{[Conditional on O-SCC.MCS and O-SCC.WeakGlue discharged,
PASS, Phase-A.]}
Let $(X_n)_{n \geq n_0}$ be a Van Hove exhaustion of SCC objects
with transported invariants $(\mathcal{T}_{X_n})$ satisfying the
five ICDC conditions of Book~III, Scholium~III.8.1.  Then the
scale-limit SCC object
$X_\infty := \varinjlim_\partial X_n$
exists and carries a coherent transported invariant
$\mathcal{T}_{X_\infty}$, with the SCC structural property
inherited from the finite stages.
\end{theorem}

\begin{proof}[Proof sketch]
Apply the ICDC Schema (Book~III, Sch.~III.8.1) with the SCC
specialization $\Theta_n \in [0,1]$ (regime-dependent).  The
five ICDC conditions (one-step inheritance, interface localization,
two-step transitivity, source-image discipline, summable residues)
are verified for the SCC transported invariant using O-SCC.MCS
(minimal carrier ensures faithful support survives the limit) and
O-SCC.WeakGlue (pairwise amalgamation ensures the direct limit
is well-defined).
\end{proof}

%% ---------------------------------------------------------------
\section{R1: SCC Branch-Projection}
\label{sec:r1}
%% ---------------------------------------------------------------

\begin{theorem}[\status{O}]\label{thm:scc-r1}
\textup{(R1: SCC Branch-Projection, Weak Form.)}
\textup{[Reserved --- no present theorem force.]}

The SCC Branch-Projection Theorem in its weak form: there exists
a canonical projection from the persistence-selected SCC sector
to the space of structural counterfactual capacity labels,
analogous to the Branch Projection Theorem (Book~V, Prop.~V.3.1)
at the SCC endpoint level.

\emph{Status:} Reserved.  This is the R1 obligation from the
retired Branch~A* governance layer.  It cannot be cited as a
discharged theorem.  It is recorded here to make the SCC
program's full obligation structure visible.  Discharge requires
MCS, UC, and TSI to be first established.
\end{theorem}

%% ---------------------------------------------------------------
\section{Final Status Proposition}
\label{sec:status}
%% ---------------------------------------------------------------

\begin{proposition}[\status{U}]\label{prop:scc-status}
\textup{(Canonical Status.)}  The SCC branch currently satisfies:
\begin{enumerate}[label=\textup{(\arabic*)}]
  \item \textbf{Lawful branch:} confirmed (Prop.~\ref{prop:scc-lawful}).
  \item \textbf{Anti-smuggling:} enforced at the primitive level
        (SR~\ref{sr:scc-no-smuggling}) and at the endpoint level
        (Thm.~\ref{thm:scc-tsifirewall}).
  \item \textbf{TSI:} structural-label form only (weak form).
        Full TSI blocked by O-SCC.MCS.  Failure code FI-TSI.
  \item \textbf{Strongest admissible endpoint label:}
        \emph{persistence-selected structural counterfactual
        capacity.}
  \item \textbf{What is NOT licensed:}
        consciousness, phenomenal experience, subjective states,
        qualia, or any stronger label.  These fail the SCC TSI
        firewall (FI-IONT, FI-IPHEN).
  \item \textbf{Status:} CERT-PROJ at most.  NOT CERT-CLOSE.
\end{enumerate}
\end{proposition}

%% ---------------------------------------------------------------
\section{Boundary of the SCC Branch Book}
%% ---------------------------------------------------------------

\textbf{What this branch book establishes.}
\begin{enumerate}
  \item The SCC branch is a lawful NFC branch (CERT-PROJ).
  \item The SCC TSI firewall (Thm.~\ref{thm:scc-tsifirewall}):
        a formal theorem governing endpoint label admissibility.
  \item The SCC anti-smuggling proposition
        (Prop.~\ref{prop:scc-antismuggle}).
  \item Branch trichotomy (Prop.~\ref{prop:scc-trichotomy}):
        the three-way classification of SCC object pairs.
  \item The four precisely named open obligations
        (Thm.~\ref{thm:scc-frontier}).
  \item The SCC-ICDC conditional theorem
        (Thm.~\ref{thm:scc-icdc}).
  \item The R1 SCC Branch-Projection theorem recorded as
        a Reserved obligation (Thm.~\ref{thm:scc-r1}).
\end{enumerate}

\textbf{What this branch book does not establish.}
\begin{enumerate}
  \item Discharge of any core obligation (WeakGlue, MCS, UC, TSI)
        as a canonical theorem.
  \item CERT-CLOSE status.
  \item Consciousness, phenomenal experience, or any stronger
        label than ``persistence-selected structural counterfactual
        capacity.''
  \item A quantum theory of consciousness or any AI-specific claim.
  \item Any empirical cognitive science import.
\end{enumerate}

%% ---------------------------------------------------------------
\section{Dependency Ledger}
%% ---------------------------------------------------------------

\renewcommand{\arraystretch}{1.4}
\noindent
\begin{tabular}{@{}p{4.0cm}p{0.8cm}p{3.6cm}p{4.2cm}@{}}
\hline
\textbf{Theorem} & \textbf{St.} & \textbf{Depends on} & \textbf{Exports to} \\
\hline
\ref{prop:scc-formation}\ Formation
  & [U] & Bk.~V Thm.~V.8.1
  & \ref{prop:scc-lawful} \\

\ref{prop:scc-label}\ Label Admissibility
  & [U] & Bk.~VII Prop.~VII.6.x
  & \ref{thm:scc-tsifirewall} \\

\ref{thm:scc-tsifirewall}\ TSI Firewall
  & [U] & \ref{prop:scc-label}; FI codes
  & Endpoint label discipline \\

\ref{prop:scc-lawful}\ SCC Branch Lawful
  & [U] & \ref{prop:scc-formation}; Bk.~IV--V
  & All subsequent SCC theorems \\

\ref{prop:scc-distinct}\ Full Distinctness
  & [U] & Terminal predicates
  & Branch catalog \\

\ref{prop:scc-trichotomy}\ Branch Trichotomy
  & [U] & $\spath$, $\scap$ definitions
  & Branch identity logic \\

\ref{prop:scc-antismuggle}\ Anti-Smuggling
  & [U] & Bk.~I SR~I.1.3; \ref{sr:scc-no-smuggling}
  & Carrier reduction; O-SCC.MCS \\

\ref{thm:scc-frontier}\ Failure Frontier
  & [U] & All obligations named
  & External scrutiny; CERT-PROJ posture \\

\ref{thm:scc-icdc}\ SCC-ICDC
  & [C] & O-SCC.MCS + WeakGlue; Bk.~III Sch.~III.8.1
  & O-SCC.UC; scale-limit SCC \\

\ref{thm:scc-r1}\ R1 (Reserved)
  & [O] & O-SCC.MCS + UC + TSI
  & Future SCC closure program \\

\ref{prop:scc-status}\ Canonical Status
  & [U] & All of SCC Branch Book
  & Program ledger \\
\hline
\end{tabular}
\renewcommand{\arraystretch}{1}

\medskip

\begin{scholium}[\status{R}]
The SCC Branch Book is frontier-honest.  The canonical program
makes no claim about consciousness and does not use consciousness
as a premise.  What it does claim is precisely delineated by the
SCC TSI firewall: persistence-selected structural counterfactual
capacity.  This is a mathematically tractable structural property
that emerges from the NFC source-descended carrier structure.
Whether this structural property bears any relation to phenomenal
consciousness is a question this book explicitly does not address,
because addressing it would require theorem content not present
in the canonical spine.  The firewall is not a philosophical
evasion; it is the canonical program's honest statement of where
its theorem-bearing force currently ends.
\end{scholium}

\end{document}
